\documentclass[paper=a4, fontsize=11pt]{jlreq} % 日本語用の標準的なクラス

% --- パッケージ群 ---
\usepackage[utf8]{inputenc}
\usepackage[T1]{fontenc}
\usepackage{amsmath, amssymb, amsthm} % 数学の基本
\usepackage{physics} % \ket{0}, \pdv{x} などが簡単に打てる
\usepackage{bm} % 太字ベクトル \bm{v}
\usepackage{graphicx}
\usepackage{hyperref} % リンク・目次用
\usepackage{cite} % 引用用

% --- 定理環境の設定 (OMCの問題風に整理できる) ---
\theoremstyle{definition}
\newtheorem{theorem}{定理}[section]
\newtheorem{definition}[theorem]{定義}
\newtheorem{example}[theorem]{例}
\newtheorem{remark}[theorem]{注釈}
\newtheorem{problem}[theorem]{問題} % 自作問題用

% --- 独自のショートカット (CFT/パンルヴェ用) ---
\newcommand{\Virexp}[2]{[L_{#1}, L_{#2}]} % ヴィラソロの交換関係
\newcommand{\taufunc}{\tau(z)} % タウ関数

\title{多様体に関する勉強ノート}
\author{松浦光佑}
\date{\today}

\begin{document}

\maketitle
\tableofcontents % 自動で目次を作成

\section{はじめに}
本ノートは、松本『多様体の基礎』に関する学習記録である。
日々の学習を「AC（Accepted）」として記録し、GitHubでの進捗可視化を目的とする。

\section{多様体の基礎}
\begin{definition}[座標近傍、局所座標系]
    % --- ここに数式 ---
    X:位相空間とする。$U\subset X$:開集合から$U'\subset\mathbb{R}^{m}$:開集合への同相写像
    \[
    \phi : U \to U'
    \]
    があるとき$(U,\phi)$を\textbf{m次元座標近傍}といい、$\phi$を$U$上の\textbf{局所座標系}という。
    \begin{itemize}
        \item[\textbf{Insight:}] 多様体上の小領域を数空間$\mathbb{R}^{m}$に対応させる「地図」であり、抽象的な図形の上で微分などの計算を可能にするための定義。地球儀の一部を平面図として扱うように、図形上の各点に「座標」という数値の番地を割り振る役割を担っている。
        \item[\textbf{Link:}] 松本6節の定義6.1/p.37を参照。
    \end{itemize}
\end{definition}

\begin{definition}[位相多様体]
    % --- ここに数式 ---
    位相空間$M$が次の条件(1)(2)を満たす時、$M$を\textbf{m次元位相多様体}という。
    \begin{itemize}
        \item[(1)] $M$はハウスドルフ空間
        \item[(2)] 任意の$p\in M$に対して$p$を含むm次元座標近傍$(U,\phi)$が存在する。
    \end{itemize}
    \begin{itemize}
        \item[\textbf{Insight:}] 「局所的には数空間 $\mathbb{R}^m$ と見なせる」という性質により、複雑な図形を慣れ親しんだ計算の土俵に載せるための定義。ハウスドルフ性は空間としての健全性（点が適切に分離できること）を保証し、全域が「地図（座標近傍）」で漏れなく覆い尽くせることを約束している。
        \item[\textbf{Link:}] 松本6節の定義6.2/p.38を参照。
    \end{itemize}
\end{definition}



% --- 今後の課題（TODOリスト） ---
\section{Next Step / To-Do}
\begin{itemize}
    \item [ ] 座標変換について学習する
\end{itemize}

\end{document}